This chapter provides a comprehensive overview of relevant concepts that are essential for this work.

\section{Microgrids}
Economic, technological, and environmental factors are driving changes in how electricity is generated and transmitted. The dominance of large centralized power plants is diminishing and giving way to smaller, more decentralized models, partly due to the decreasing advantages of traditional economies of scale \cite{MicrogridConcept}. The concept of a microgrid was first presented by Lasseter \cite{Microgrid} as early as 2002, describing a paradigm to define the operation of distributed generation. The core idea is to perceive such generators and associated loads as a system of its own to allow for local control, thereby providing a higher local reliability and efficiency.

The CIGRÉ Working Group C6.22 \textit{Microgrid Evolution Roadmap}  defines the following components as the foundational elements of a microgrid \cite{MicrogridKeyComponents}:
\begin{itemize}
    \item[] \textbf{Generators} encompass various energy sources suitable for microgrids, including fossil or biomass-fired small-scale combined heat and power, photovoltaic modules, wind turbines, etc.
    \item[] \textbf{Storage Devices} include diverse technologies such as electrical, electrochemical, mechanical, and heat storage.
    \item[] \textbf{Controlled loads} refer to all kinds of electrical loads within a microgrid that can be managed or controlled to optimize energy usage and system performance.
\end{itemize}

In addition to that, a microgrid can operate in two different modes: grid-connected and islanded \cite{MicrogridConcept, MicrogridKeyComponents, MicrogridComprehensiveStudy}. In the grid-connected mode, the microgrid functions as a part of the larger utility grid, while drawing power from the grid as needed, enabling it to benefit from the main grids stability and reliability. During certain phases like grid outages, the microgrid is capable of a seamless transition into islanded mode where it disconnects from the grid and operates autonomously, relying solely on its own distributed energy resources to meet local energy demands. Effective control and coordination mechanisms are thereby essential for smooth transitions between these two modes, to optimize energy management and ensure reliable operation under varying conditions.